\section{Grading and Assessment}
\textit{Reference $|$ Course Text: Chapters 11, 13}

Grading in Capstone Design is inherently more complex than in a lecture course. You are assessing engineering judgment, teamwork, and professional conduct---not just deliverable completeness. This section explains the grading structure, rubric philosophy, and common calibration challenges.

\subsection{How Final Grades Are Determined}

Final grades are determined by the CF using a combination of inputs:

\begin{enumerate}[nosep,leftmargin=1.5em]
\item \textbf{PA input.} Your assessment of team performance across sprints, design reviews, and deliverables. This is the most heavily weighted qualitative input.
\item \textbf{Deliverable scores.} Rubric-based scores on reports, presentations, and assignments (graded on Canvas).
\item \textbf{Feedback Loop.} Peer evaluations and Individual Performance Evaluations that differentiate individual contributions within a team.
\item \textbf{CF observation.} The CF's direct observation at design reviews and through course-level data.
\end{enumerate}

The CF synthesizes these inputs into final grades. PAs do not assign final grades, but your input is the primary basis for the qualitative assessment. Take this responsibility seriously.

\begin{warnbox}
Do not promise or imply specific grades to students. You may describe quality levels (``this is strong work,'' ``this needs significant improvement''), but grade assignment is the CF's responsibility. If a student presses you for a grade commitment, redirect them to the syllabus and the CF.
\end{warnbox}

\subsection{Rubric Philosophy}

The grading rubrics evaluate two dimensions:

\begin{enumerate}[nosep,leftmargin=1.5em]
\item \textbf{Engineering quality.} Does the work demonstrate sound engineering judgment? Are trade-offs analyzed with data? Are assumptions stated and justified? Are conclusions supported by evidence?
\item \textbf{Professional execution.} Is the deliverable complete, well-organized, properly formatted, and free of errors? Would you be comfortable sending this to a client?
\end{enumerate}

\begin{bestpractice}
When reviewing deliverables, ask yourself: ``Would I accept this from a junior engineer on my team?'' This framing calibrates expectations appropriately---better than a student report, but not yet at the level of a licensed professional engineer.
\end{bestpractice}

\subsection{Assessing Engineering Calculation Packages}

The engineering calculation package is where teams demonstrate analytical rigor. An A-level calculation package includes:

\begin{itemize}[nosep,leftmargin=1.5em]
\item \textbf{Stated assumptions.} Every assumption is listed, and the team can justify why each is reasonable.
\item \textbf{Method justification.} The team explains why they chose this analytical method (FEA, hand calc, empirical correlation, etc.) and acknowledges its limitations.
\item \textbf{Uncertainty analysis.} Results include uncertainty bounds or sensitivity analysis. A single point estimate without uncertainty is incomplete.
\item \textbf{Validation.} Where possible, the analysis is compared against experimental data, published results, or an independent calculation method.
\item \textbf{Design connection.} The conclusion explicitly connects the analysis result to a requirement or design decision. ``The factor of safety is 2.3, which meets the requirement of $\geq$2.0'' is useful. ``The maximum stress is 145~MPa'' without context is not.
\end{itemize}

\subsection{Differentiating Individual Contributions}

Team grades are the starting point, but individual grades may differ. Use the following evidence to differentiate:

\begin{itemize}[nosep,leftmargin=1.5em]
\item \textbf{Feedback Loop peer evaluations.} Cross-reference peer scores with your own observations. Consistent low scores from multiple peers are a reliable signal.
\item \textbf{Sprint review participation.} Who answers questions with depth? Who defers to others on their own assigned topics?
\item \textbf{Deliverable ownership.} Review version histories and commit logs. Who is writing, who is editing, and who is absent?
\item \textbf{Meeting attendance and engagement.} Chronic absence or passive attendance is a differentiator.
\item \textbf{PA Individual Performance Evaluation.} Your direct assessment of each student, completed at scheduled intervals throughout the semester.
\end{itemize}

\begin{tipbox}
Keep brief notes after each sprint review: one sentence per student summarizing their contribution that sprint. These notes are invaluable at the end of the semester when you need to differentiate grades and cannot rely on memory alone.
\end{tipbox}

\subsection{Handling Grade Disputes}

Grade disputes occasionally arise, particularly when individual grades differ from team grades. Follow this process:

\begin{enumerate}[nosep,leftmargin=1.5em]
\item \textbf{Listen first.} Let the student explain their concern fully before responding.
\item \textbf{Refer to evidence.} Point to specific deliverable feedback, peer evaluation scores, and your observation notes. Avoid subjective characterizations without supporting data.
\item \textbf{Escalate to the CF.} If the student is not satisfied with your explanation, direct them to the CF. The CF makes final grade determinations and handles formal appeals.
\item \textbf{Do not change grades unilaterally.} Any grade changes go through the CF.
\end{enumerate}

\subsection{The Feedback Loop: Interpretation and Application}

The Feedback Loop peer evaluation system generates quantitative scores (1--5 scale) and qualitative comments from teammates. Here is how to interpret and apply the data:

\textbf{Interpreting scores:}
\begin{itemize}[nosep,leftmargin=1.5em]
\item \textbf{4.0--5.0:} Strong contributor. Peer perception aligns with high performance.
\item \textbf{3.0--3.9:} Adequate contributor. May have specific areas for improvement but is pulling their weight.
\item \textbf{Below 3.0:} Consistent underperformance perceived by peers. Cross-reference with your own observations before acting.
\end{itemize}

\textbf{Applying scores:}
\begin{itemize}[nosep,leftmargin=1.5em]
\item A single low evaluation may reflect interpersonal friction, not performance. Look for \emph{patterns} across multiple evaluation cycles.
\item If peer scores and your observations diverge significantly, investigate. The team may be shielding a struggling member, or there may be a personality conflict distorting scores.
\item Share aggregate trends (not individual scores) with the student: ``Your peer evaluations suggest that your teammates would like to see more proactive communication.''
\end{itemize}

\begin{warnbox}
A pattern of low Feedback Loop scores (below 3.0) combined with PA observation of non-participation can result in individual grade reduction, PIP initiation, or course failure---regardless of team deliverable scores. Document your observations throughout the semester.
\end{warnbox}


% ═══════════════════════════════════════
